% ---------
%  Compile with "pdflatex hw1".
% --------
%!TEX TS-program = pdflatex
%!TEX encoding = UTF-8 Unicode

% Template borrowed from Jeff Erickson.

\documentclass[11pt]{article}
\usepackage{jeffe,handout,graphicx}
\usepackage[utf8]{inputenc}		% Allow some non-ASCII Unicode in source
\usepackage{xcolor}
\usepackage{mdframed}
\usepackage{arydshln}

% =========================================================
%   Define common stuff for solution headers
% =========================================================
\Class{CS 6319.001}
\Semester{Spring 2024}
\Authors{1}
\AuthorOne{Christian Parker}{cxp220034}
%\Section{}

% =========================================================
\begin{document}
\HomeworkHeader{2}{1}% homework number, problem number
% ---------------------------------------------------------
\begin{enumerate}[(a)]\itemsep0pt
  \item
    Suppose you are given a triangulation with \(v\) vertices, \(h\) of which lie on the outer face.
    Use Euler's formula to derive the number of triangles \(t\) and the number of edges \(e\) in
    the triangulation as functions of \(v\) and h.
    
\begin{solution}

  There's always one more face than there are triangles because of the external face.

  \(f = t + 1\)

  Euler's formula:

  \(v - e + f = 2\)

  \(v - e + t + 1 = 2\)

  \(v - e + t = 1\)

  \(2v - 2e + 2t = 2\)

  Total number of edges incident to all triangles:

  \(3t = 2e - h\)

  \(2e = 3t + h\)

  Replace \(2e\) with \(3t + h\) in the equation \(2v - 2e + 2t = 2\):

  \(2v - 3t - h + 2t = 2\)

  \(2v - t - h = 2\)

  \(t = 2v - h - 2\)

  Plug \(t\) back into \(v - e + t = 1\):

  \(v - e + 2v - h - 2 = 1\)

  \(3v - e - h = 3\)

  \(e = 3v - h - 3\)
\end{solution}

  \item
    Suppose you are given a quandrangulation with \(v\) vertices, \(h\) of which lie on the outer
    face. Use Euler’s formula to derive the number of quadrangles \(q\) and the number of
    edges \(e\) in the quadrangulation as functions of \(v\) and h.

\begin{solution}
  Total number of edges incident to all quadrangles:

  \(4q = 2e - h\)

  \(2e = 4q + h\)

  Euler's formula (skipping steps shown in (a)):

  \(v - e + q = 1\)

  \(2v - 2e + 2q = 2\)

  Replace \(2e\) with \(4q + h\) in the equation \(2v - 2e + 2q = 2\):

  \(2v - 4q - h + 2q = 2\)

  \(2v - 2q - h = 2\)

  \(2q = 2v - h - 2\)

  \(q = v - \frac{h}{2} - 1\)

  Plug \(q\) back into \(v - e + q = 1\):

  \(v - e + v - \frac{h}{2} - 1 = 1\)

  \(2v - e - \frac{h}{2} = 2\)

  \(e = 2v - \frac{h}{2} - 2\)

\end{solution}

  \item
    Explain why your answer to part (b) implies that no quadrangulation exists where
    the number of vertices on the outer face is odd.

\begin{solution}
  If \(h\) was odd, then \(\frac{h}{2}\) would not be an integer, therefore both the number of quadrangles and the number of edges would not be integers.
  We couldn't have a partial quadrangle or edge though so \(h\) must not be odd.
\end{solution}

\end{enumerate}

% ---------------------------------------------------------
\HomeworkHeader{2}{2}% homework number, problem number
% ---------------------------------------------------------
\begin{enumerate}[(a)]\itemsep0pt
  \item
    Compute the line equations of the two "crossing tangents", that is, the lines
    \(l_1: y = a_1 x - b_1\) and \(l_2: y = a_2 x - b_2\) that are both supporting lines for the convex
    hulls of each of \(P\) and \(Q\) such that \(P\) lies below \(l_1\) and above \(l_2\) and the reverse holds
    for \(Q\). (Note that you are not given the hulls, just the point sets.) Your algorithm
    should run in \(O(n + m)\) time.
    
\begin{solution}
  
\end{solution}

  \item
    You have a cannon in \(\mathbb{R}^2\). It has three controls labeled "a", "b", and "c". A projectile
    shot from this cannon travels along the parabolic arc \(y = a + bx - c x^2\). You are
    asked to determine whether it is possible to adjust the controls so that the projectile
    travels above a set of \(n\) building tops, represented by a point set \(P = \{p_1,..., p_n\}\) and
    beneath a set of \(m\) floating balloons, represented by a point set \(Q = \{q_1,...,q_m\}\). Your
    algorithm should run in time \(O(n + m)\). (Do not worry about the cannon's location.
    Just determine if there is some parabola along which the projectile can travel.)

\begin{solution}

\end{solution}

  \item
    You are given a set of \(n\) halfplanes \(H = \{h1,...,hn\}\). Each halfplane \(h_i\) is given as a
    pair \((a_i , b_i)\) and consists of all the points \((x, y)\) such that \(y \leq a_i x - b_i\). Compute the
    axis-parallel square of the largest side length lying in \(h_1 \cap ... \cap h_n\) whose lower edge
    lies on the x-axis. If no such square exists, your algorithm should indicate this.

\begin{solution}
  
\end{solution}

\end{enumerate}

% ---------------------------------------------------------
\HomeworkHeader{2}{3}% homework number, problem number
% ---------------------------------------------------------
\begin{enumerate}[(a)]\itemsep0pt
  \item
    Suppose we are given \(n\) points and create a subdivision as described above. Show
    that the resulting subdivision has size \(O(n)\) (counting vertices, edges, and faces).
    
\begin{solution}
  Each time we add a point to the subdivision, there are three possible cases:
  \begin{enumerate}
    \item
      The point is inside an existing face. In this case, we shoot bullets left, right, and up. Each of these creates a new vertex when it hits an edge.
      Add in the point itself and that's 4 new vertices. Each bullet creates a new edge so 3 new edges. There's now a face below the horizontal line and
      to the left and right of the vertical line. This adds 2 faces.

      So in total we add 4 vertices, 3 edges, and 2 faces.
    \item
      The point is on a horizontal edge. In this case, we split the edge into two edges and shoot a bullet up. The face is split into two faces. One new vertex is created
      at the end of the bullet path and a vertex is added for the new point.

      So in total we add 2 vertices, 2 edges, and 1 face.
    \item
      The point is on a vertical edge. In this case, we split the edge into two edges and shoot a bullet left and right. The left and right faces are split into two faces each.
      One new vertex is created at the end of each bullet path and a vertex is added for the new point.

      So we add 3 vertices, 3 edges, and 2 faces.
  \end{enumerate}
  In the worst case, we add 4 vertices, 3 edges, and 2 faces for each point. So the size of the subdivision is \((4 + 3 + 2) \cdot n = 9n = O(n)\).
\end{solution}

  \item
    Prove that for any \(1 \leq i \leq n\), the expected
    number of new boxes created during the ith addition is \(O(1)\).

\begin{solution} Pulling heavily from Mount's notes because this is basically the same as the proof that the expected number of new trapezoids created during
  the \(i\)th addition to a trapezoidal map is \(O(1)\).

  \textbf{Lemma:} Consider the randomized incremental construction of a subdivision given points \(P = \{p_1,...,p_n\}\), and let \(k_i\) denote the
  number of new boxes created when the \(i\)th point is added. Then \(E[k_i] = O(1)\), where the expectation is taken over all possible permutations of the
  points as the insertion orders.

  \textbf{Proof:} Let \(S_i\) denote the subdivision resulting from the insertion of the \(i\)th point. Because we are averaging over all permutations, among the \(i\) points
  in \(S_i\), each point has an equal probability \(\frac{1}{i}\) of being the last one added. For each of the points \(p\), we want to count the number of boxes
  that would have been created if \(p\) had been the last point added (The worst case scenario for \(p\)).

  We say that a box \(b\) depends on a point \(p\) if \(p\) would have caused \(b\) to be created had \(p\) been the last point inserted. We want to count the number of boxes
  that depend on each point and then compute the average over all points. If we let \(\delta(b,p)=\) 1 if \(b\) depends on \(p\) and 0 otherwise, then the expected value is:

  \(E[k_i] = \sum_{p \in P_i}\) (probability that \(p\) is last) \(\cdot\) (number of boxes that depend on \(p\))

  Where \(P_i\) is the set of points we've considered so far. We can rewrite this as:

  \begin{align*}
    E[k_i] &= \sum_{p \in P_i} \frac{1}{i} \left( \sum_{b \in S_i} \delta(b,p) \right) \\
    &= \frac{1}{i} \sum_{p \in P_i} \sum_{b \in S_i} \delta(b,p) \\
    &= \frac{1}{i} \sum_{b \in S_i} \sum_{p \in P_i} \delta(b,p)
  \end{align*}

  Each box is bounded by lines from at most 4 points. It may be less than 4 if the box is bounded by the bounding rectangle instead. The top and bottom of each box are formed by bullet paths going left
  or right from two points (or by the boundary rectangle of course). The left and right are formed by bullet paths going up from two points.
  
  Therefore, \(\sum_{p \in P_i} \delta(b,p) \leq 4\).

  Since \(S_i\) consists of at most \(2i + 1\) boxes (we start with 1 box and, as stated in part (a), add at most 2 boxes for each point), we have:

  \begin{align*}
    E[k_i] \leq \frac{1}{i} \sum_{b \in S_i} 4 = \frac{4}{i} |S_i| \leq \frac{4}{i}(2i + 1) = O(1)
  \end{align*}

\end{solution}

\end{enumerate}

\end{document}